%\documentclass[11pt,a4paper,twoside]{article}

\usepackage[utf8]{inputenc}

\usepackage[T1,T2A]{fontenc}

\usepackage[english.ancient,russian.ancient]{babel}

\usepackage{graphicx}

\usepackage{array}

\usepackage[doi]{samgtu-bib}

\usepackage{samgtu}

\usepackage{samgtu-name}

\usepackage{mathtools}

\mathtoolsset{showonlyrefs}

\usepackage{desclist}

\usepackage{euscript}

\usepackage{mathrsfs} 

\usepackage{multicol}

\usepackage{mathrsfs}

\usepackage{cite}

\usepackage{cmap}

\usepackage{qrcode}

\allowdisplaybreaks[4]

\usepackage[nodayofweek]{datetime}

\usepackage[thinlines]{easybmat}



\renewcommand{\JournalYear}{2018}

\renewcommand{\JournalVolume}{22}

\renewcommand{\JournalNumber}{4}



\newdate{JournalDateSign}{29}{12}{2018}% Дата подписания Вестника в печать

\renewcommand{\JournalDateSign}{\displaydate{JournalDateSign}}

\newdate{JournalDatePrint}{16}{01}{2019}% Дата выхода Вестника из печати

\renewcommand{\JournalDatePrint}{\displaydate{JournalDatePrint}}

%\renewcommand{\JournalUPL}{16.56} % Усл. печ. л.

%\renewcommand{\JournalUIL}{16.53} % Уч.-изд. л.

\renewcommand{\JournalUPL}{15.24} % Усл. печ. л.

\renewcommand{\JournalUIL}{15.21} % Уч.-изд. л.

\renewcommand{\JournalRegNumber}{220/18} % Регистрационный номер

\renewcommand{\JournalOrderNumber}{2} % Номер заказа



%\renewcommand{\JournalRMonth}{Март} % Месяц выхода Вестника

%\renewcommand{\JournalEMonth}{March} % Месяц выхода Вестника

%\renewcommand{\JournalRMonth}{Июнь} % Месяц выхода Вестника

%\renewcommand{\JournalEMonth}{June} % Месяц выхода Вестника

%\renewcommand{\JournalRMonth}{Сентябрь} % Месяц выхода Вестника

%\renewcommand{\JournalEMonth}{September} % Месяц выхода Вестника

\renewcommand{\JournalRMonth}{Декабрь} % Месяц выхода Вестника

\renewcommand{\JournalEMonth}{December} % Месяц выхода Вестника



\newcommand{\totop}[1]{\raisebox{6pt}[1pt][2pt]{#1}}

\newcommand{\tottop}[1]{\raisebox{12pt}[1pt][2pt]{#1}} 

\newcommand{\totttop}[1]{\raisebox{18pt}[1pt][2pt]{#1}} 

\newcommand{\tottttop}[1]{\raisebox{24pt}[1pt][2pt]{#1}} 

\newcommand{\totttttop}[1]{\raisebox{30pt}[1pt][2pt]{#1}} 

\newcommand{\tobot}[1]{\raisebox{-6pt}[1pt][2pt]{#1}} 

\newcommand{\tobbot}[1]{\raisebox{-12pt}[1pt][2pt]{#1}} 

\newcommand{\tobbbot}[1]{\raisebox{-18pt}[1pt][2pt]{#1}}



\graphicspath{{./00PIC/}}

%

%\usepackage{samgtu-name}

%

\vsgtu{1574} %registration number

\FirstPage{651}

\LastPage{664}

%

%\pdfcrop

%\draft

\ResearchArticle

%

%

%\begin{document}



\hfill \qrcode[hyperlink,height=.8in]{http://doi.org/10.14498/vsgtu1574}

\vspace{-.8in}





\udc{517.956.326}

\msc{35L25, 35L80}



\rutitle{Краевая задача для гиперболического уравнения третьего порядка с~вырождением порядка внутри области}

\rutitleshort{Краевая задача для гиперболического уравнения третьего порядка\dots}



\entitle{A boundary value problem for a third order hyperbolic equation with degeneration of order inside the domain}

\entitleshort{A boundary value problem for a third order hyperbolic equation\dots}



\ruauthor{Р.~Х.~Макаова}{М\,а\,к\,а\,о\,в\,а~Р.~Х.}

\enauthor{R.~Kh.~Makaova}{M\,a\,k\,a\,o\,v\,a~R.~Kh.}



\ruaffil{\runiipma.}



\enaffil{\enniipma.}



\ruauthorsdetails{1}{% Количество авторов

\fullauthor{\rupid{82995}{Рузанна Хасанбиевна Макаова}}

\CAuthor % Автор-корреспондент

\orcid{0000-0003-4095-2332}

%\regalia{кандидат физико-математических наук, доцент} 

\position[n]{младший научный сотрудник} 

\dept{отд. уравнений смешанного типа} 

\email[n]{makaova.ruzanna@mail.ru}

}



\enauthorsdetails{1}{

\fullauthor{\enpid{82995}{Ruzanna Kh. Makaova}}

\CAuthor % Сorresponding Author

\orcid{0000-0003-4095-2332} % ORCID ID

%\regalia{Cand. Phys. \& Math. Sci., Associate Professor} 

\position[n]{Junior Researcher} 

\dept{Dept. of Mixed Type Equations}

\email{makaova.ruzanna@mail.ru}

}



\rukeywords{краевая задача, гиперболическое уравнение третьего порядка, уравнение Аллера}

\enkeywords{boundary-value problem, third order hyperbolic equation, Hallaire equation}



\ruabstract{Исследуется краевая задача для гиперболического уравнения третьего порядка с~вырождением типа внутри смешанной области. Рассматриваемое уравнение в~положительной части области совпадет с~уравнением Аллера, которое является уравнением псевдопараболического типа. А~в~отрицательной части области "--- с~вырождающимся гиперболическим уравнением первого рода, частным случаем которого является уравнение Бицадзе"--~Лыкова. 

Доказана теорема существования и~единственности решения. Единственность решения задачи доказана с~помощью метода Трикоми. Из функциональных соотношений, принесенных на линию вырождения порядка из положительной и~отрицательной частей области, приходим к~уравнению Вольтерра второго рода типа свертки относительно следа производной искомого решения. 

Путем применения метода интегрального преобразования Лапласа  решение интегрального уравнения находится в~явном виде. 

Далее решение исследуемой задачи выписывается в~явном виде как решение второй краевой задачи для уравнения Аллера в~положительной части области и~как решение задачи Коши для вырождающегося гиперболического уравнения первого рода в~отрицательной части области.}



\enabstract{In this paper we study the boundary value problem for a degenerating third order equation of hyperbolic type in a mixed domain. The equation under consideration in the positive part of the domain coincides with the Hallaire equation, which is a pseudoparabolic type equation. Moreover, in the negative part of the domain it coincides with a degenerating hyperbolic equation of the first kind, the particular case of the Bitsadze--Lykov equation. The existence and uniqueness theorem for the solution is proved. The uniqueness of the solution to the problem is proved with the Tricomi method. Using the functional relationships of the positive and negative parts of the domain on the degeneration line, we arrive at the convolution type Volterra integral equation of the 2nd kind with respect to the desired solution by a derivative trace. With the Laplace transform method, we obtain the solution of the integral equation in its explicit form. At last, the solution to the problem under study is written out explicitly as the solution of the second boundary-value problem in the positive part of the domain for the Hallaire equation and as the solution to the Cauchy problem in the negative part of the domain for a degenerate hyperbolic equation of the first kind.}



\newdate{ruzanaa}{27}{10}{2017}

\date{\displaydate{ruzanaa}}

\newdate{ruzanab}{11}{12}{2017}

\revisiondate{\displaydate{ruzanab}}

\newdate{ruzanac}{18}{12}{2017}

\accepteddate{\displaydate{ruzanac}}



\newdate{ruzanae}{28}{12}{2017}

\renewcommand{\JournalDataOnline}{\displaydate{ruzanae}}



%\ForPeerReview

\makerutitle



\newpage



\Section[n]{Введение}

В евклидовой плоскости точек $(x,y)$ рассматривается уравнение вида

\begin{equation}

\label{Makaova1}

0=\left\{

\begin{array}{ll}

u_y-au_{xx}-bu_{xxy}, & y>0, \\

(-y)^m u_{xx}-u_{yy}-c(-y)^{({m-2})/{2}}u_x, & y<0, \\

\end{array}

\right.  

\end{equation}

где $a$, $b$, $m$  "--- заданные положительные числа;

$ |c |\leq m/2$;

$u=u(x,y)$ "--- искомая действительная функция независимых переменных $x$ и $y$.



Уравнение \eqref{Makaova1} при $y>0$  совпадает с  уравнением Аллера~\cite{Hallaire_1}:

\begin{equation}

\label{Makaova2} 

\frac{\partial u}{\partial y}

=\frac{\partial }{\partial x}

\Bigl(

a\frac{\partial u}{\partial x}

+b\frac{\partial^2 u}{\partial x \partial y}

\Bigr), 

\end{equation}

а при $y<0$  "--- с  вырождающимся гиперболическим уравнением первого рода~\cite{Smirnov_2}:

\begin{equation}

\label{Makaova3} 

(-y)^m u_{xx}-u_{yy}-c(-y)^{({m-2})/{2}}u_x=0.

\end{equation}



Уравнение \eqref{Makaova2} относят к уравнениям псевдопарабролического типа \cite{Showalter_3}, 

хотя оно является уравнением гиперболического типа. 

При определенных допущениях уравнение~\eqref{Makaova2} описывает движение влаги в~капиллярно"=пористых средах~\cite{Chudnovsky_4} и  его решение <<интерпретируется как влажность почвы с~коэффициентом диффузии~$a$ и~коэффициентом влагопроводности $b$ в~точке $x$ почвенного слоя $0\leq x \leq r$ в момент времени $y \in [0,T]$>>~\cite{Nakhushev_5}.  

Изучению краевых задач для уравнений третьего порядка псевдопараболического типа, в частности для уравнения Аллера,  посвящены такие работы, как~[\citen{Barenblatt_6}--\citen{Kozhanov_12}]. 



Уравнение \eqref{Makaova3} является  уравнением гиперболического типа и~при $m=2$ его называют уравнением 

Бицадзе"--~Лыкова~[\citen{Nakhushev_13}--\citen{Repin_15}].

При $c=0$ уравнение~\eqref{Makaova3} переходит в~уравнение Геллерстедта, которое находит применение к~отыскании оптимальной формы плотины прорези~\cite{Nakhushev_16}. 

В~работах \cite{Kalmenov_17, Kalmenov_18} были изучены первая и~вторая  задачи Дарбу для уравнения \eqref{Makaova3}.

Исследованию различных краевых задач для вырождающихся гиперболических уравнений посвящены работы 

[\citen{Smirnov_2}, \citen{Pulkina_19}--\citen{Balkizov_22}].



\smallskip





\Section[n]{Постановка задачи и полученный результат} Пусть

$\Omega^+=\{ (x,y):$ \linebreak 



\vspace{-3mm}





\begin{minipage}{0.35\textwidth}

	\begin{center}

		\includegraphics[scale=0.75]{makaovapic}

		%\centerline{Рис.~1.}

	\end{center}

\end{minipage}

\hfill

\begin{minipage}{0.60\textwidth}

	\hspace{1em} %  абзацный отступ внутри minipage



$ 0<x<r, 0<y<T \}$. 

 Через $\Omega^-$ обозначим область, ограниченную характеристиками уравнения \eqref{Makaova3}:  

\[A_0C: x-\frac{2}{m+2}(-y)^{({m+2})/{2}}=0,

\]

\[A_rC: x+\frac{2}{m+2}(-y)^{({m+2})/{2}}=r,

\]

выходящими из точек $A_0=(0, 0)$, $A_r=(r, 0)$ и пересекающимися в точке 

$$C=\bigl(r/2, - [(m+2)r/4]^{2/(m+2)} \bigr),$$ и отрезком $A_0A_r=\{(x,0): 0<x<r \}$. 

Пусть $B_0=(0,T)$, 

$B_r=(r,T)$ и $A_0B_0=\{(0,y):0<y<T\}$, 

$A_rB_r=\{(r,y):0<y<T\}$; $\Omega=\Omega^+ \cup \Omega^- \cup (A_0A_r)$.	

\end{minipage}









\begin{definition}

Регулярным в области $\Omega$ решением уравнения \eqref{Makaova1} назовем функцию $u=u(x,y)$ такую, что $u \in C( \overline{\Omega}) \cap C^1(\Omega) \cap C^2(\Omega^-)$, $u_{xx}$, $u_{xxy} \in C(\Omega^+)$, $u_{xx}(x,0)$, $u_{xxy}(x,0) \in L_1(A_0A_r)$,  и удовлетворяющую уравнению~\eqref{Makaova1}.

\end{definition} 



\smallskip



Исследуется следующая 

\smallskip

\begin{task}

Найти регулярное в области $\Omega$ решение уравнения \eqref{Makaova1} из класса $u_{x} \in C(\overline{A_0B_0} \cup \overline{A_rB_r}), $

удовлетворяющее следующим условиям{\rm:}

\begin{gather}

\label{Makaova4} 

u_x(0,y)=\nu(y), \quad  0\leq y\leq T, 

\\

\label{Makaova5} 

u_x(r,y)=\nu_r(y), \quad  0\leq y\leq T, 

\\

\label{Makaova6} 

u|_{CA_r}=h_r(x), \quad   r/2 \leq x\leq r, 

\end{gather}

где  $\nu(y)$, $\nu_r(y) \in C^1[0,T]$, $h_r(x)\in C^3[0,r]$.

\end{task}



\smallskip





Имеет место следующая 

\smallskip



\begin{theorem} Регулярное решение задачи  \eqref{Makaova4}--\eqref{Makaova6} для уравнения \eqref{Makaova1}

существует и оно единственно\rm.

\end{theorem} 



\smallskip



\begin{proof}

{\bf Единственность решения. }

Положим 

\begin{gather}

\label{Makaova7} 

u(x,0)=\varphi(x),  \quad  0\leq x\leq r,  

\\

\label{Makaova8} 

u_y(x,0)=\psi(x),  \quad   0< x\leq r.

\end{gather}



Из \eqref{Makaova2}, переходя к пределу при $y \to +  0$ с учетом \eqref{Makaova7} и \eqref{Makaova8}, находим функциональное соотношение между функциями $ \varphi(x)$ и $\psi(x)$,

принесенное из области $\Omega^+$ на линию $y=0$, в виде

\begin{equation}

\label{Makaova9} 

\psi(x)-a\varphi''(x)-b\psi''(x)=0.

\end{equation}

Проинтегрировав равенство \eqref{Makaova9} дважды  в пределах от $x$ до $r$, получим

\begin{equation}

\label{Makaova10} 

a\varphi(x)+b\psi(x)-\int _r^x(x-\xi) \psi(\xi) d \xi=

a\varphi(r)+b\psi(r)-(r-x)\left[a\varphi'(r)+b\psi'(r)\right].

\end{equation}



Пользуясь формулой нахождения производной по

заданному направлению \cite[п.~184]{Fikhtengolz_23}, запишем следующую формулу для функции $h_r(x)$:

\begin{equation}

\label{Makaova11}

h_r'(x)=u_y(x,y)y_x +u_x(x,y), \quad  \forall (x,y)\in CA_r.

\end{equation}



Для любых точек $(x,y) \in  CA_r$ из \eqref{Makaova11} получаем 

\begin{equation}

\label{Makaova12}

u_y(x,y)=

\left[h_r'(x)-u_x(x,y)\right] \Bigl[\frac{m+2}{2}(r-x)\Bigr]^{{m}/({m+2})}.

\end{equation}

Так как $u_y(x,0) \in L_1(A_0A_r)$,  из формулы \eqref{Makaova12} верно, 

что ${u_y(r,0)=\psi(r)=0}$. 

Тогда с~учетом условий согласования

$$\varphi(r)=h_r(r),\quad \varphi'(r)=\nu_r(0), \quad \psi'(r)=\nu_r'(0)$$ 

равенство \eqref{Makaova10}

можно переписать в виде

\begin{equation}

\label{Makaova13} 

\varphi(x)= \frac{1}{a} D_{rx}^{-2} \psi(\xi)-\frac{b}{a} \psi(x)+

h_r(r)-(r-x)\Bigl[ \nu_r(0)+\frac{b}{a}\nu_r'(0)\Bigr].

\end{equation}







Через $D^\gamma_{rx}$  обозначим оператор дробного интегро"=дифференцирования в~смысле

Римана"--~Лиувилля по переменной  $x$ порядка $\gamma $ с~началом в~точке $r$ и с~концом в~точке $x$, определяемый следующим образом \cite{Nakhushev_16}:

\[

 D_{rx}^{\gamma}  v(\xi)=\frac{\mathop{\rm sign}(x-r)}{\Gamma(-\gamma)} \int _{r}^{x} 

 \frac{v(\xi) d \xi}{|x-\xi|^{\gamma+1}}, \quad   \gamma<0;

\]

\[

 D_{rx}^{\gamma}  v(\xi)=v(x), \quad   \gamma=0;

\]

\[

 D_{rx}^{\gamma}  v(\xi)=\mathop{\rm sign}^n(x-r) \frac{d^n}{dx^n} 

 D_{rx}^{\gamma-n}  v(\xi), \quad   \gamma \in ]n-1, n], \;  n=1, 2,\dots,

\]

%\[

%D_{0x}^\gamma v(\xi,y) =\frac{\partial^n}{\partial x^n}

%D_{0x}^{\gamma-n}  v(\xi,y), \ \ \ n-1<\gamma\leq n,\ \ n=1,2, ...;

%\]

%\[

%D_{0x}^{\gamma} v(\xi,y) =  v(x,y), \ \ \ \gamma=0, 

%\]

где 

$\Gamma(z)$  "--- гамма-функция Эйлера.



Рассмотрим интеграл

\begin{equation}

\label{Makaova14} 

I=\int _0^r u(x,0) u_y(x,0) dx=\int _0^r \varphi(x) \psi(x) dx.

\end{equation} 

Из \eqref{Makaova13} и \eqref{Makaova14} при однородных краевых условиях 

\eqref{Makaova4}--\eqref{Makaova6}  получим

\begin{equation}

\label{Makaova15} 

I=\frac{1}{a} \int _0^r 

\psi(x)

D_{rx}^{-2} \psi(\xi)d x

-\frac{b}{a}\|\psi(x)\|^2_0,

\end{equation}

где

$\displaystyle \|f\|^2_0=\int _0^r f^2(x)dx$  "---  норма в~$L_2[0,r]$.



Интеграл, стоящий в правой части равенства \eqref{Makaova15}, перепишем  в виде

\begin{multline}

(\psi, D^{-2}_{rx}\psi)_0 =

\int _0^r \psi(x)  \int _x^r (\xi-x) \psi(\xi) d\xi d x=

\\

=

\int _0^r \psi(\xi)d\xi  \int _0^\xi (\xi-x) \psi(x) d x  =

(\psi, D^{-2}_{0x}\psi)_0,

\end{multline}

откуда находим

\begin{equation}

\label{Makaova16} 

   (\psi, D^{-2}_{0x}\psi)_0 =

\int _0^r (r-x) \psi(x) d x \int _0^r  \psi(x) dx -

\int _0^r \biggl(\int _0^x  \psi(\xi) d \xi \biggr)^2 dx,

\end{equation}

где $\displaystyle (f,g)_0=\int _0^r f(x) g(x)dx$ "--- скалярное произведение в пространстве $L_2[0,r]$.





Далее проинтегрируем уравнение \eqref{Makaova2} по переменной $x$ в пределах от $0$ до~$r$. 

С~учетом  \eqref{Makaova4} и \eqref{Makaova5} имеем

\[

\int _0^r u_y(x,y) dx=\int _0^r \frac{\partial}{\partial x}

\Bigl[ a\frac{\partial u}{\partial x}+b\frac{\partial^2 u}{\partial x \partial y}\Bigr]dx

=a\nu_r(y)+b\nu_r'(y)-a\nu(y)-b\nu'(y).

\]

Отсюда, когда условия \eqref{Makaova4} и \eqref{Makaova5} однородные, находим

\begin{equation}

\label{Makaova17}

\int _0^r u_y(x,y) dx=0, \quad  0\leq y\leq T.

\end{equation}

Следовательно, из \eqref{Makaova16} с учетом условия \eqref{Makaova17} получаем оценку

\[

\left( \psi, D_{rx}^{-2}\psi \right)_0=\left( \psi, D_{0x}^{-2}\psi \right)_0=-

\biggl\|\int _0^x \psi(\xi) d\xi \biggr\|^2_0 \leq 0,

\]

при использовании которой из \eqref{Makaova15} следует неравенство

\begin{equation}

\label{Makaova18}

I=

-\frac{b}{a}\|\psi(x)\|^2_0 -\frac{1}{a} \biggl\| \int _0^x \psi(\xi) d \xi \biggr\|^2_0\leq 0.

\end{equation}



Для того чтобы получить соотношение между $\varphi(x)$ и $\psi (x)$, принесенное из области $\Omega^-$ на линию $y=0$, выпишем решение задачи Коши \eqref{Makaova7}, \eqref{Makaova8} для уравнения \eqref{Makaova3} \cite{Smirnov_2}:

\begin{multline} 

u(x,y)=\frac{1}{B(\alpha, \beta)}\int _0^1 

\varphi \Bigl[ x+\frac{2}{m+2}(-y)^{({m+2})/{2}}(2t-1)\Bigr] t^{\beta-1} (1-t)^{\alpha-1}dt+

\\

\label{Makaova19}

+\frac{y}{B(1-\alpha,1- \beta)} \int _0^1 

\psi \Bigl[ x+\frac{2}{m+2}(-y)^{({m+2})/{2}}(2t-1)\Bigr] \times \\

\times t^{-\alpha} (1-t)^{-\beta}dt, \quad  |c|<\frac{m}{2};

\end{multline}

\begin{multline} 

u(x,y)=\varphi \Bigl[ x+\frac{2}{m+2}(-y)^{({m+2})/{2}}\Bigr]+

\\

\label{Makaova21} 

+\frac{2y}{m+2}

\int _0^1 

\psi \Bigl[ 

x+\frac{2}{m+2}(-y)^{({m+2})/{2}}(2t-1)

\Bigr]

(1-t)^{-\beta}dt, \quad  c=\frac{m}{2};

\end{multline}

\begin{multline} 

u(x,y)=\varphi \Bigl[ x-\frac{2}{m+2}(-y)^{({m+2})/{2}}

\Bigr]+

\\

\label{Makaova21} 

+\frac{2y}{m+2}

\int _0^1 

\psi \Bigl[ x+\frac{2}{m+2}(-y)^{({m+2})/{2}}(1-2t)\Bigr]

(1-t)^{-\alpha}dt, \quad  c=-\frac{m}{2},

\end{multline} 

где 

\[

\alpha=\frac{m-2c}{2(m+2)}, \quad \beta=\frac{m+2c}{2(m+2)},

\]

а $B(\gamma, z)$  "--- бета-функция.  



Учитывая условие \eqref{Makaova6}, перепишем \eqref{Makaova19}--\eqref{Makaova21} в виде

\begin{multline} 

h_r\Bigl(\frac{r+x}{2} \Bigr)=

\frac{(r-x)^{1-\alpha-\beta} \Gamma(\beta)}{B(\alpha, \beta)} D^{-\beta}_{rx}

\bigl[ (r-\xi)^{\alpha-1} \varphi(\xi) \bigr]-

\\

\label{Makaova22} 

-\frac{\Gamma(1-\alpha)}{B(1-\alpha, 1-\beta)} 

\Bigl( \frac{m+2}{4} \Bigr)^{{2}/({m+2}) }

 D^{\alpha-1}_{rx}

\bigl[ (r-\xi)^{-\beta} \psi(\xi) \bigr],  \quad   |c|<\frac{m}{2};

\end{multline}

\begin{gather}

\label{Makaova24} 

h_r\Bigl(\frac{r+x}{2} \Bigr)=

\varphi(r)-\frac{1}{2} 

\Bigl( \frac{m+2}{4} \Bigr)^{-\beta} 

 D^{-1}_{rx}   \bigl[(r-\xi)^{-\beta} \psi(\xi)\bigr],  \quad  c=\frac{m}{2};

\\

\label{Makaova24} 

h_r\Bigl(\frac{r+x}{2} \Bigr)=

\varphi(x)-\frac{\Gamma(1-\alpha)}{2} 

\Bigl( \frac{m+2}{4} \Bigr)^{-\alpha} 

 D^{\alpha-1}_{rx}  \psi(\xi) ,  \quad  c=-\frac{m}{2}.

\end{gather}



Для любой функции $v(x) \in L_1[0,r]$ справедливы следующие свойства   дробного интегро"=дифференцирования с~одинаковыми началами \cite{Nakhushev_13}:

\[

D^{\alpha_1}_{rx}D^{-\beta_1}_{r\xi} v(s)=D^{\alpha_1-\beta_1}_{rx}v(\xi),\quad 0<\alpha_1 \leq \beta_1;

\]

\[

D^{\alpha_1}_{rx} | \xi-r |^{\alpha_1+\beta_1}D^{\beta_1}_{r\xi} v(s)=

|x-r |^{\beta_1}D^{\alpha_1+\beta_1}_{rx} |\xi-r |^{\alpha_1} v(\xi),\quad 

\beta_1<0, \;  0\leq \alpha_1<1,

\]

с применением которых из \eqref{Makaova22}--\eqref{Makaova24}

функция $\psi(x)$ выражается следующим образом:

\begin{gather}

\label{Makaova25} 

 \psi(x)=

 \frac{d_2}{d_1}  D^{1-\alpha-\beta}_{rx} \varphi(\xi)-

 \frac{(r-x)^{\beta} }{d_1} D^{1-\alpha}_{rx} h_r \Bigl(\frac{r+\xi}{2} \Bigr),   \quad |c|<\frac{m}{2};

\\

\label{Makaova27} 

 \psi(x)=

- \frac{(r-x)^{\beta} }{2d_3}  h_r'\Bigl(\frac{r+x}{2} \Bigr),   \quad  c=\frac{m}{2};

\\

\label{Makaova27} 

 \psi(x)=

 \frac{1}{d_4}  D^{1-\alpha}_{rx} \varphi(\xi)-

 \frac{1 }{d_4} D^{1-\alpha}_{rx} h_r\Bigl(\frac{r+\xi}{2} \Bigr),   \quad  c=-\frac{m}{2},

 \end{gather}

 где

\[d_1=\Bigl( \frac{m+2}{4} \Bigr)^{{2}/({m+2}) } \frac{\Gamma(2-\alpha-\beta)}{\Gamma(1-\beta)} , \quad 

d_2=\frac{\Gamma(\alpha+\beta)}{\Gamma(\alpha)} ,

\]

\[

d_3=\frac{1}{2} \Bigl( \frac{m+2}{4} \Bigr)^{-\beta } , \quad 

d_4=\frac{1}{2} \Bigl( \frac{m+2}{4} \Bigr)^{-\alpha } \Gamma(1-\alpha).

\]

Подставляя найденные значения  \eqref{Makaova25}--\eqref{Makaova27} функции 

$\psi(x)$  при $h_r(x)=0$, из \eqref{Makaova14} имеем:



\begin{gather}

\label{Makaova28} 

I=\frac{d_2}{d_1} \int _0^r \varphi(x) 

D^{1-\alpha-\beta}_{rx} \varphi(\xi) dx, \quad 

 |c|<\frac{m}{2};

\\

\label{Makaova29} 

I= 0,  \quad c=\frac{m}{2};

\\

\label{Makaova30} 

I=\frac{1}{d_4} \int _0^r \varphi(x)

   D^{1-\alpha}_{rx} \varphi(\xi) dx,    \quad  c=-\frac{m}{2}.

 \end{gather}

 

Известно \cite{Balkizov_24}, что для любой абсолютно непрерывной на  $[0,r]$ 

функции $v(x)$ при $v(r)=0$ справедливо неравенство 

 \begin{equation}

\label{Makaova31}  v(x)D_{rx}^\gamma v(\xi) \geq \frac{1}{2}D_{rx}^\gamma v^2(\xi), \quad  0<\gamma \leq 1.

 \end{equation}

Тогда с учетом  \eqref{Makaova31} и представления

 \[

 D_{rx}^\gamma v^2(\xi)=-\frac{1}{\Gamma(1-\gamma)}

 \int _x^r \frac{2v(\xi)v'(\xi)}{(\xi-x)^\gamma}d\xi

 \]

из \eqref{Makaova28} и \eqref{Makaova30} имеем 

\begin{multline}

\label{Makaova32} 

I \geq 

\frac{d_2}{2d_1}

\int _0^r D_{rx}^{1-\alpha-\beta} \varphi^2(\xi)

dx= \\ =

\frac{d_2}{2d_1\Gamma(\alpha+\beta)}

\int _0^r x^{\alpha+\beta-1} \varphi^2(x)

dx \geq 0, \quad  |c|<\frac{m}{2};

 \end{multline}

 

 \begin{equation}

\label{Makaova33} 

I\geq 

\frac{1}{2d_4}

\int _0^r D_{rx}^{1-\alpha} \varphi^2(\xi)

dx=

\frac{1}{2d_4\Gamma(\alpha)}

\int _0^r x^{\alpha-1} \varphi^2(x)

dx \geq 0, \quad    c=-\frac{m}{2}.

 \end{equation}

Таким образом, из  \eqref{Makaova32} и \eqref{Makaova33} получаем, что $I\geq 0$.

С другой стороны, из  \eqref{Makaova18} имеем, что $I \leq 0$. 

Следовательно, $I=0$ при $-m/2 \leq c < m/2$ , которое имеет место тогда и только тогда, когда $\varphi(x)=0$. 

Тогда, как видно из \eqref{Makaova25} и~\eqref{Makaova27}, функция $\psi(x)=0$.

При $c=m/2$ из   \eqref{Makaova18}  и \eqref{Makaova29} следует, что $\psi(x)=0$,

и~тогда  из \eqref{Makaova13} получаем, что и $\varphi(x)=0$.  

При этом из формул \eqref{Makaova19}--\eqref{Makaova21} имеем $u(x,y) \equiv 0$ в области $\Omega^-$.





В области $\Omega^+$

задача для уравнения \eqref{Makaova1} совпадает со второй краевой задачей для уравнения Аллера, решение которого выписывается 

в виде \cite{Makaova_25}:

\begin{multline}

u(x, y)=\int _{0}^{r}

\left [ \varphi(\xi)-b \varphi''(\xi)

 \right ]G(x, y; \xi,0)d\xi+

\\

\label{Makaova34}

+\int _{0}^{y}  

G(x, y; r,\eta) \left [a \nu_r(\eta)+b \nu_r'(\eta) \right ]d\eta- \\

-

\int _{0}^{y}

G(x, y; 0, \eta) \left [a\nu(\eta)+b \nu'(\eta) 

\right ] d\eta,

 \end{multline}

 где

\begin{multline}

G(x,y; \xi, \eta)

=\frac{1}{r}  

+\frac{2}{r} 

\sum_{n=1}^{\infty} 

\frac{1}{1+b\mu_n}

\exp \Bigl({ -\frac{a \mu_n}{1+b\mu_n} (y-\eta) }\Bigr) \times \\ \times

\cos \left ( \sqrt{ \mu_n} \xi \right ) \cos \left (\sqrt{ \mu_n} x  \right ), 

\ \mu_n =\left (\frac{\pi n}{r} \right )^2 . 

 \end{multline}



При однородных условиях   \eqref{Makaova4}, \eqref{Makaova5} и \eqref{Makaova7} из представления \eqref{Makaova34} находим, что  $u(x,y) \equiv 0$ в области $\Omega^+$. Таким образом, однородная задача соответствующая задаче \eqref{Makaova4}--\eqref{Makaova6} для уравнения \eqref{Makaova1}, имеет только тривиальное решение, откуда следует единственность регулярного решения исследуемой задачи.



\smallskip



{\bf Существование решения.} 

Перейдем к доказательству существования решения исследуемой задачи.

Из \eqref{Makaova22}--\eqref{Makaova24} верно, что

 \begin{multline}

\label{Makaova35} 

\varphi(x)=\frac{d_1}{d_2}

D_{rx}^{\alpha+\beta-1}\psi(\xi)+ \\ 

+

\frac{(r-x)^{1-\alpha}}{d_2}

D_{rx}^\beta

\Bigl[(r- \xi)^{\alpha+\beta-1} h_r \Bigl(\frac{r+\xi}{2} \Bigr)

\Bigr], \quad  |c|<\frac{m}{2};

\end{multline}



\begin{gather}

\label{Makaova37} 

\varphi(r)=d_3D_{rx}^{-1} 

\Bigl[

(r-\xi)^{-\beta}\psi(\xi)\Bigr]+h_r\Bigl(\frac{r+x}{2} \Bigr)

,

\quad  c=\frac{m}{2};

\\

\label{Makaova37} 

\varphi(x)=d_4D_{rx}^{\alpha-1}\psi(\xi)+

h_r\Bigl(\frac{r+x}{2} \Bigr), \quad  c=-\frac{m}{2}.

\end{gather}

Исключая функцию $\varphi(x)$ из \eqref{Makaova13} и соотношений \eqref{Makaova35}--\eqref{Makaova37}, 

относительно  функции $\psi(x)$ получим следующие равенства:

 \begin{gather}

\label{Makaova38} 

\psi(x)+\frac{ad_1}{bd_2}D^{\alpha+\beta-1}_{rx}\psi(\xi)-\frac{1}{b}D^{-2}_{rx}\psi(\xi)=f_\beta(x), 

\quad  |c|<\frac{m}{2};

\\

\label{Makaova39} 

\psi(x)=-\frac{(r-x)^\beta}{2d_3} h_r'\Bigl( \frac{r+x}{2} \Bigr), \quad  c=\frac{m}{2};

\\

\label{Makaova40} 

\psi(x)+\frac{a}{bd_4}D^{\alpha-1}_{rx}\psi(\xi)-\frac{1}{b}D^{-2}_{rx}\psi(\xi)=f_0(x), \quad   c=-\frac{m}{2},

 \end{gather}

где 

\begin{multline*}

f_\beta(x)=\frac{a}{b}h_r(r)-(r-x)

\Bigl[\frac{a}{b} \nu_r(0)+\nu_r'(0)\Bigr]-

\\

-\frac{a(r-x)^{1-\alpha}}{bd_2}D^\beta_{rx} 

\Bigl[

(r-\xi)^{\alpha+\beta-1} h_r\Bigl( \frac{r+\xi}{2} \Bigr)

\Bigr].

\end{multline*}



С учетом равенства

 \[

D^\gamma_{r \ r-x}\psi(\xi)=

D^\gamma_{0 x}\psi(r-\xi), \quad  \gamma<0

\]

после замены $x$   на  $(r-x)$, когда $\psi(r-x)=q(x)$ и   $f_\beta(r-x)=w(x)$,

 из  \eqref{Makaova38} получим

  \begin{equation}

\label{Makaova41} 

q(x)+\frac{ad_1}{bd_2}D^{\alpha+\beta-1}_{0x}q(\xi)-\frac{1}{b}D^{-2}_{0x}q(\xi)=w(x).

 \end{equation} 



Уравнение \eqref{Makaova41} представляет собой уравнение Вольтерра второго рода

типа свертки. 

Следовательно, при  $w(x) \in L_1[0,r]$  оно имеет единственное решение.



В силу представления оператора дробного интегрирования перепишем уравнение \eqref{Makaova41} в виде 

 \begin{equation}

\label{Makaova42} 

q(x)+\frac{ad_1}{bd_2} \bigl(q(x)* x^{-\alpha-\beta}\bigr)-\frac{1}{b}

\bigl(q(x)*x\bigr)=\bigl(x*W(x)\bigr). 

 \end{equation} 

Здесь $ \displaystyle \bigl (f(x)*g(x)\bigr)=\int _{0}^{x}f(x-\xi)g(\xi) d\xi $  "--- свертка функций 

$f(x)$ и $g(x)$; $W(x)\hm =D^2_{0x}w(\xi)$.



Пусть $Q(p)$ и $W^*(p)$ являются изображениями функций $q(x)$ и 

$W(x)$ по Лапласу:

$$

q(x)\risingdotseq Q(p),\quad  W(x) \risingdotseq W^*(p).$$

Применяя преобразование Лапласа, 

из \eqref{Makaova42} находим

\begin{equation}

\label{Makaova43} 

Q(p)=\frac{ W^*(p)}{\Delta(p)p^{2}},

 \end{equation} 

 где 

\[

\Delta(p)=1+\frac{ad_1}{bd_2p^{1-\alpha-\beta}}-\frac{1}{bp^{2}}.

\]

Так как при достаточно больших значениях $p$  имеет место равенство 

\[

\int _{0}^{\infty}e^{-\Delta (p)t}dt=\frac{1}{\Delta (p)},

\]

 равенство \eqref{Makaova43} можно переписать в виде

\begin{equation}

\label{Makaova44} 

Q(p)=\int _{0}^{\infty}\frac{e^{-\Delta(p)t}}

{p^{2}} W^*(p)dt.

 \end{equation} 



Найдем  обратное преобразование  Лапласа к уравнению \eqref{Makaova44}. Для этого  воспользуемся известной формулой $x^{k-1} \phi(\gamma, k;zx^\gamma) 

\ \risingdotseq p^{-k}e^{zp^{-\gamma}}$ \cite{Wright_26}:  

\begin{equation}

\label{Makaova45} 

q(x)=\int _{0}^{\infty}e^{-t} 

\Bigl[

W(x)*\phi \Bigl(1-\alpha-\beta;1; -\frac{ad_1}{bd_2} tx^{1-\alpha-\beta}\Bigr)* \\

\phi \Bigl(2;1;\frac{1}{b} tx^2 \Bigr)

\Bigr] dt,

\end{equation}

где

\[

\phi(\gamma, k;t) =\sum^\infty_{n=0}

\frac{t^n}{n! \Gamma(n\gamma+k)}

\]

"--- функция  Райта.







Воспользовавшись  обозначениями, введенными в работе \cite{Pskhu_27}:

\begin{gather}

G_n^k \equiv G_n^k(x; \lambda_1, \dots, \lambda_n;\gamma_1, \dots, \gamma_n)=

\int _{0}^{\infty}e^{-t} S_n^k(x; \lambda_1t, \dots, \lambda_nt;\gamma_1, \dots, \gamma_n)dt,

\\

S_n^k(x; \lambda_1t, \dots, \lambda_nt;\gamma_1, \dots, \gamma_n)=

(h_1 \ast h_2 \ast \cdots \ast h_n)(x),

\\

h_i=h_i(x)=x^{k_i-1} \phi(\gamma_i, k_i; \lambda_itx^{\gamma_i}), \quad 

i= 1,2, \dots, n;\quad   k=\sum _{i=1}^n k_i,

\end{gather}

решение \eqref{Makaova45}  можно переписать в виде



\begin{multline}

\psi(r-x)=q(x)= 

\int _{0}^{\infty}e^{-t} \Bigl[W(x)* S_2^2 \Bigl( x; -\frac{ad_1}{bd_2}t, \frac{1}{b}t;1-\alpha-\beta,2 \Bigr)\Bigr]

dt=

\\

= \int _{0}^{x} W(\xi)  G_2^2 \Bigl( x-\xi; -\frac{ad_1}{bd_2}, \frac{1}{b};1-\alpha-\beta,2 \Bigr) d\xi.

\end{multline}

Таким образом, переобозначив в последнем равенстве $(r-x)$ через $x$ и выполнив

подстановку $t=r-\xi$ под знаком интеграла, получим

\begin{multline}

\psi(x)= \int _{0}^{r-x} w^{''}(\xi)  G_2^2

\Bigl( r-x-\xi; -\frac{ad_1}{bd_2}, \frac{1}{b};1-\alpha-\beta,2 \Bigr) d\xi=

\\

\label{Makaova46}

=\int _{x}^{r} f_\beta^{''}(t) G_2^2

\Bigl( t-x; -\frac{ad_1}{bd_2}, \frac{1}{b};1-\alpha-\beta,2 \Bigr) dt.

\end{multline}

















При $\beta=0$ из \eqref{Makaova46} получаем решение  уравнения 

 \eqref{Makaova40} в виде

\begin{equation}

\label{Makaova47} 

\psi(x)=\int _{r}^{x} f_0^{''}(t)  G_2^2

\Bigl( t-x; -\frac{ad_1}{b}, \frac{1}{b};1-\alpha,2 \Bigr) dt.

\end{equation}





Из \eqref{Makaova13}, \eqref{Makaova35}  и \eqref{Makaova37} при подстановке формул \eqref{Makaova39}, \eqref{Makaova46} и \eqref{Makaova47}

 функция $\varphi(x)$ находится однозначно. Следовательно, после нахождения функций $\varphi(x)$ и $\psi (x)$ решение задачи \eqref{Makaova4}--\eqref{Makaova6} для уравнения \eqref{Makaova2} 

в области $\Omega^+ $ находится по формуле \eqref{Makaova34}, а в области 

$\Omega^-$ решение задачи Коши для уравнения \eqref{Makaova3} выписывается по одной из формул \eqref{Makaova19}--\eqref{Makaova21}.

 \end{proof} 







\bigskip



%\AdditionalContent{Конкурирующие интересы}{Укажите какие конфликты интересов имеются.} 

\AdditionalContent{Конкурирующие интересы}{Конкурирующих интересов не имею.} 

%\AdditionalContent{Авторский вклад и ответственность}{Укажите степень авторской ответственности каждого автора.} 

\AdditionalContent{Авторская ответственность}{Я несу полную ответственность за предоставление окончательной версии рукописи в печать. Окончательная версия рукописи мною одобрена.} 

%\AdditionalContent{Авторский вклад и ответственность}{Все авторы принимали участие в разработке концепции \mbox{статьи} и в написании рукописи. Авторы несут полную ответственность за предоставление окончательной рукописи в печать. Окончательная версия рукописи была одобрена всеми авторами.} 

%\AdditionalContent{Финансирование}{Укажите источник финансирования работы.}

\AdditionalContent{Финансирование}{Исследование выполнялось без финансирования.}

%\AdditionalContent{Финансирование}{Работа выполнена при поддержке РФФИ (проект № xx–xx–xxxxx_a).}

%\AdditionalContent{Благодарность}{Выразите свою благодарность.}

%\AdditionalContent{Благодарность}{Авторы благодарны рецензенту за тщательное прочтение статьи и~ценные предложения и комментарии.}







\begin{thebibliography}{99}



\Bibitem{Hallaire_1}

\by Hallaire~M.

\paper Potential efficace de l'eau dans le sol en r\'egime de dess\`echement 

\inbook \href{http://hydrologie.org/redbooks/062.htm}{Assembl\'ee g\'en\'erale de Berkeley=General Assembly of Berkeley}\rm, \nofrills

\procinfo August 1963

\bookinfo Publ. no. 62

\yr 1963

\publaddr Gentbrugge

\pages 114--122 

\elink \url{http://hydrologie.org/redbooks/a062/iahs_062_0114.pdf}



\RBibitem{Smirnov_2}

\by Смирнов~М.~М.

\book Вырождающиеся гиперболические уравнения

\yr 1977

\publ Вышэйшая школа

\publaddr Минск

\totalpages 158







\Bibitem{Showalter_3}

\by Showalter R.E., Ting T.W.

\paper Pseudoparabolic partial differential equations

\jour SIAM J. Math. Anal.

\yr 1970

\vol 1

\issue 1

\pages 1--26

\crossref{10.1137/0501001}



\RBibitem{Chudnovsky_4}

\by Чудновский~А.~Ф.

\book Теплофизика почв

\yr 1976

\publ Наука

\publaddr М.

\totalpages 352







\RBibitem{Nakhushev_5}

\by Нахушев~А.~М. 

\paper Об одном классе нагруженных уравнений в частных производных дробного порядка

\jour Доклады Адыгской (Черкесской) Международной академии наук

\yr 2012

\vol 14

\issue 1

\pages 51--57

\elib{https://elibrary.ru/item.asp?id=17749735}





\RBibitem{Barenblatt_6}

\by Баренблатт~Г.~И., Желтов~И.~П., Кочина~И.~Н.

\paper Об основных представлениях теории фильтрации однородных жидкостей в трещиноватых породах

\jour ПММ 

\vol 24

\issue 5

\yr 1960

\pages 852–864

%\by Barenblatt G.I., Zheltov Iu.P., Kochina I.N.

%\paper Basic concepts in the theory of seepage of homogeneous liquids in fissured rocks [strata]

%\jour PMM

%\yr 1960

%\vol 24

%\issue 5

%\pages 852--864







\Bibitem{Coleman_7}

\crossref{10.1007/BF00282277}

\by  Coleman~B.~D., Duffin~R.~J., Mizel~V.~J.

\paper IInstability, uniqueness, and nonexistence theorems for the equation $u_t=u_{xx}-u_{xtx}$ on a strip

\jour  Arch. Rat. Mech. Anal.

\yr 1965

\vol 19

\issue 2

\pages 100--116



\Bibitem{Colton_8}

\by Colton D.

\paper Pseudoparabolic equations in one space variable

\jour J.~Differ. Equations

\yr 1972

\vol 12

\issue 3

\pages 559--565

\crossref{10.1016/0022-0396(72)90025-3}

		

\RBibitem{Yangarber_9}

\by Янгарбер~В.~А.

\paper О смешанной задаче для модифицированного уравнения влагопереноса

\jour ПМТФ

\yr 1967

\issue 1

\pages 91--96



\RBibitem{Shkhanukov_10}

\by Шхануков~М.~Х.

\paper О~некоторых краевых задачах для уравнения третьего порядка, возникающих при моделировании

фильтрации жидкости в~пористых средах

\jour Дифференц. уравнения

\yr 1982

\vol 18

\issue 4

\pages 689--699

\mathnet{http://mi.mathnet.ru/de4523}

\mathscinet{http://www.ams.org/mathscinet-getitem?mr=658436}

\zmath{https://zbmath.org/?q=an:0559.76088}





\RBibitem{Vodakhova_11}

\by Водахова~В.~А.

\paper Краевая задача с~нелокальным условием А.\,М.~Нахушева для одного псевдопараболического уравнения влагопереноса

\jour Дифференц. уравнения

\yr 1982

\vol 18

\issue 2

\pages 280--285

\mathnet{http://mi.mathnet.ru/de4442}

\mathscinet{http://www.ams.org/mathscinet-getitem?mr=649672}

\zmath{https://zbmath.org/?q=an:0486.35045}



\RBibitem{Kozhanov_12}

\by Кожанов~А.~И.

\paper Об одной нелокальной краевой задаче с~переменными коэффициентами для уравнений

теплопроводности и Аллера

\jour Дифференц. уравнения

\yr 2004

\vol 40

\issue 6

\pages 763--774

\mathnet{http://mi.mathnet.ru/de11086}

\mathscinet{http://www.ams.org/mathscinet-getitem?mr=2162443}

%\transl

%\jour Differ. Equ.

%\yr 2004

%\vol 40

%\issue 6

%\pages 815--826

%\crossref{10.1023/B:DIEQ.0000046860.84156.f0}



\RBibitem{Nakhushev_13}

\by Нахушев~А.~М.

\book Уравнения математической биологии

\yr 1995

\publ Высшая школа

\publaddr М.

\totalpages 301







\RBibitem{Repin_14}

\by О.~А.~Репин

\paper Аналог задачи Нахушева для уравнения Бицадзе--Лыкова

\jour Дифференц. уравнения

\yr 2002

\vol 38

\issue 10

\pages 1412--1417

\mathnet{http://mi.mathnet.ru/de10721}

\mathscinet{http://www.ams.org/mathscinet-getitem?mr=2014239}

%\transl

%\jour Differ. Equ.

%\yr 2002

%\vol 38

%\issue 10

%\pages 1503--1509

%\crossref{10.1023/A:1022339217281}









\RBibitem{Repin_15}

\by Репин~О.~А., Кумыкова~С.~К.

\paper Нелокальная задача для уравнения Бицадзе--Лыкова

\jour Изв. вузов. Матем.

\yr 2010

\issue 3

\pages 28--35

\mathnet{http://mi.mathnet.ru/ivm6709}

\mathscinet{http://www.ams.org/mathscinet-getitem?mr=2778322}

%\transl

%\jour Russian Math. (Iz. VUZ)

%\yr 2010

%\vol 54

%\issue 3

%\pages 24--30

%\crossref{10.3103/S1066369X10030059}

%\scopus{http://www.scopus.com/record/display.url?origin=inward&eid=2-s2.0-78649533261}





\RBibitem{Nakhushev_16}

\by Нахушев~А.~М.

\book Дробное исчисление и его применение

\yr 2003

\publ Физматлит

\publaddr М.

\totalpages 272









\RBibitem{Kalmenov_17}

\by Кальменов~Т.~Ш.

\paper Критерий единственности решения задачи Дарбу для одного вырождающегося

гиперболического уравнения

\jour Дифференц. уравнения

\yr 1971

\vol 7

\issue 1

\pages 178--181

\mathnet{http://mi.mathnet.ru/de1193}

\mathscinet{http://www.ams.org/mathscinet-getitem?mr=283396}

\zmath{https://zbmath.org/?q=an:0212.44402}







\RBibitem{Kalmenov_18}

\by Кальменов~Т.~Ш.

\paper О~задаче Дарбу для одного вырождающегося уравнения

\jour Дифференц. уравнения

\yr 1974

\vol 10

\issue 1

\pages 59--68

\mathnet{http://mi.mathnet.ru/de2071}

\mathscinet{http://www.ams.org/mathscinet-getitem?mr=340843}

\zmath{https://zbmath.org/?q=an:0269.35066}





\RBibitem{Pulkina_19}

\by Л.~С.~Пулькина

\paper Об~одной нелокальной задаче для вырождающегося гиперболического уравнения

\jour Матем. заметки

\yr 1992

\vol 51

\issue 3

\pages 91--96

\mathnet{http://mi.mathnet.ru/mz4503}

\mathscinet{http://www.ams.org/mathscinet-getitem?mr=1172230}

\zmath{https://zbmath.org/?q=an:0812.35085}

%\transl

%\jour Math. Notes

%\yr 1992

%\vol 51

%\issue 3

%\pages 286--290

%\crossref{10.1007/BF01206393}



\RBibitem{Repin_20}

\by О.~А.~Репин, С.~К.~Кумыкова

\paper Задача со смещением для вырождающегося внутри области гиперболического уравнения

\jour Вестн. Сам. гос. техн. ун-та. Сер. Физ.-мат. науки

\yr 2014

\issue 1(34)

\pages 37--47

\mathnet{http://mi.mathnet.ru/vsgtu1280}

\crossref{10.14498/vsgtu1280}

\elib{http://elibrary.ru/item.asp?id=22813958}



\RBibitem{Balkizov_21}

\by Балкизов~Ж.~А.

\paper Первая краевая задача для вырождающегося внутри области гиперболического уравнения

\jour Владикавк. матем. журн.

\yr 2016

\vol 18

\issue 2

\pages 19--30

\mathnet{http://mi.mathnet.ru/vmj577}



\RBibitem{Balkizov_22}

\by Балкизов~Ж.~А.

\paper Краевая задача для вырождающегося внутри области гиперболического уравнения

\jour Известия вузов. Северо-Кавказский регион. Серия: Естественные науки

\yr 2016

\issue 1

\pages 5--10

\elib{https://elibrary.ru/item.asp?id=25792370}

\crossref{10.18522/0321-3005-2016-1-5-10}









\RBibitem{Fikhtengolz_23}

\by Фихтенгольц~Г.~М.

\book Курс дифференциального и интегрального исчисления

\bookvol 1

\yr 2003

\publ Физматлит

\publaddr М.

\totalpages 680





\RBibitem{Balkizov_24}

\by Ж.~А.~Балкизов

\paper Первая краевая задача для уравнения параболо-гиперболического типа третьего порядка с вырождением типа и порядка в области гиперболичности

\jour Уфимск. матем. журн.

\yr 2017

\vol 9

\issue 2

\pages 25--39

\mathnet{http://mi.mathnet.ru/ufa373}

\elib{http://elibrary.ru/item.asp?id=29419139}

%\transl

%\jour Ufa Math. Journal

%\yr 2017

%\vol 9

%\issue 2

%\pages 25--39

%\crossref{10.13108/2017-9-2-25}

%\scopus{http://www.scopus.com/record/display.url?origin=inward&eid=2-s2.0-85023623624}







\RBibitem{Makaova_25}

\by Макаова Р.Х.

\paper Вторая краевая задача для обобщенного уравнения Аллера с дробной производной Римана--Лиувилля

\jour Доклады Адыгской (Черкесской) Международной академии наук

\yr 2015

\vol 17

\issue 3

\pages 35--38

\elib{https://elibrary.ru/item.asp?id=25117227}







\Bibitem{Wright_26}

\by   Wright~E.~M.

\paper The generalized Bessel function of order greater than one

\jour Quart. J. Math. Oxford Ser.

\yr 1940

\vol os-11

\issue 1

\pages 36--48

\crossref{10.1093/qmath/os-11.1.36}







\RBibitem{Pskhu_27}

\by Псху~А.~В.

\paper Начальная задача для линейного обыкновенного дифференциального уравнения дробного порядка

\jour Матем. сб.

\yr 2011

\vol 202

\issue 4

\pages 111--122

\mathnet{http://mi.mathnet.ru/msb7645}

\crossref{10.4213/sm7645}

\mathscinet{http://www.ams.org/mathscinet-getitem?mr=2830238}

\zmath{https://zbmath.org/?q=an:1226.34005}

\adsnasa{http://adsabs.harvard.edu/cgi-bin/bib_query?2011SbMat.202..571P}

\elib{http://elibrary.ru/item.asp?id=19066272}

%\transl

%\by Pskhu~A.~V.

%\paper Initial-value problem for a~linear ordinary differential equation of noninteger order

%\jour Sb. Math.

%\yr 2011

%\vol 202

%\issue 4

%\pages 571--582

%\crossref{10.1070/SM2011v202n04ABEH004156}

%\isi{http://gateway.isiknowledge.com/gateway/Gateway.cgi?GWVersion=2&SrcApp=PARTNER_APP&SrcAuth=LinksAMR&DestLinkType=FullRecord&DestApp=ALL_WOS&KeyUT=000292829300005}

%\scopus{http://www.scopus.com/record/display.url?origin=inward&eid=2-s2.0-79959815598}





\end{thebibliography}





\newpage



\makeentitle





%\AdditionalContent{Competing interests}{Specify what conflicts of interest are available.} 

\AdditionalContent{Competing interests}{I have no competing interests.} 

%\AdditionalContent{Competing interests}{We have ... (or) no competing interests.} 

%\AdditionalContent{Authors' contributions and responsibilities}{What is the author's responsibility for each author?}

\AdditionalContent{Author's Responsibilities}{I take full responsibility for submitting the final manuscript in print. I approved the final version of the manuscript.} 

%\AdditionalContent{Authors' contributions and responsibilities}{Each author has participated in the article concept development and in the manuscript writing. The authors are absolutely responsible for submitting the final manuscript in print. Each author has approved the final version of manuscript.} 



%\AdditionalContent{Funding}{Indicate the source(s) of your funding.}

\AdditionalContent{Funding}{The research has not had any funding.}

%\AdditionalContent{Funding}{This work was supported by the Russian Foundation for Basic Research (project no. xx–xx–xxxxx_a).}



%\AdditionalContent{Acknowledgments}{Express your gratitude (acknowledgments) here.}

%\AdditionalContent{Acknowledgments}{The authors are grateful to the referee for careful reading of the paper and valuable suggestions and comments.}







\begin{thebibliography}{99}



\Bibitem{Hallaire_1:en}

\by Hallaire~M.

\paper Potential efficace de l'eau dans le sol en r\'egime de dess\`echement 

\inbook \href{http://hydrologie.org/redbooks/062.htm}{Assembl\'ee g\'en\'erale de Berkeley=General Assembly of Berkeley}\rm, \nofrills

\procinfo August 1963

\bookinfo Publ. no. 62

\yr 1963

\publaddr Gentbrugge

\pages 114--122 

\elink \url{http://hydrologie.org/redbooks/a062/iahs_062_0114.pdf}



\Bibitem{Smirnov_2:en}

\lang In Russian

\by Smirnov~M.~M.

\book Vyrozhdaiushchiesia giperbolicheskie uravneniia \rm [Degenerate hyperbolic equations]

\yr 1977

\publ Vysheishaia shkola

\publaddr Minsk

\totalpages 158









\Bibitem{Showalter_3:en}

\by Showalter R.E., Ting T.W.

\paper Pseudoparabolic partial differential equations

\jour SIAM J. Math. Anal.

\yr 1970

\vol 1

\issue 1

\pages 1--26

\crossref{10.1137/0501001}



\Bibitem{Chudnovsky_4:en}

\lang In Russian

\by Chudnovskii~A.~F.

\book Teplofizika pochv \rm [Thermal Physics of Soils]

\yr 1976

\publ Nauka

\publaddr Moscow

\totalpages 352









\Bibitem{Nakhushev_5:en}

\lang In Russian

\by Nakhushev~A.~M. 

\paper On a class of loaded partial differential equations of fractional order

\jour Doklady Adygskoi (Cherkesskoi) Mezhdunarodnoi akademii nauk

\yr 2012

\vol 14

\issue 1

\pages 51--57









\Bibitem{Barenblatt_6:en}

\by Barenblatt~G.~I., Zheltov~Iu.~P., Kochina~I.~N.

\paper Basic concepts in the theory of seepage of homogeneous liquids in fissured rocks [strata]

\jour J. Appl. Math. Mech.

\yr 1960

\vol 24

\issue 5

\pages 1286--1303

\crossref{10.1016/0021-8928(60)90107-6}







\Bibitem{Coleman_7}

\crossref{10.1007/BF00282277:en}

\by  Coleman~B.~D., Duffin~R.~J., Mizel~V.~J.

\paper IInstability, uniqueness, and nonexistence theorems for the equation $u_t=u_{xx}-u_{xtx}$ on a strip

\jour  Arch. Rat. Mech. Anal.

\yr 1965

\vol 19

\issue 2

\pages 100--116



\Bibitem{Colton_8:en}

\by Colton D.

\paper Pseudoparabolic equations in one space variable

\jour J.~Differ. Equations

\yr 1972

\vol 12

\issue 3

\pages 559--565

\crossref{10.1016/0022-0396(72)90025-3}

		

\Bibitem{Yangarber_9:en}

\paper The mixed problem for a modified moisture-transfer equation

\by Yangarber~V.~A.

\jour J. Appl. Mech. Tech. Phys.

\yr 1967

\vol 8

\issue 1

\pages 62–64 

\crossref{10.1007/BF00913245}



\Bibitem{Shkhanukov_10:en}

\lang In Russian

\by Shkhanukov~M.~Kh.

\paper Some boundary value problems for a third-order equation that arise in the modeling of the filtration of a fluid in porous media

\jour Differ. Uravn.

\yr 1982

\vol 18

\issue 4

\pages 689--699





\Bibitem{Vodakhova_11:en}

\lang In Russian

\by Vogahova~V.~A.

\paper A boundary value problem with A.~M.~Nakhushev's nonlocal condition for a pseudoparabolic equation of moisture transfer

\jour Differ. Uravn.

\yr 1982

\vol 18

\issue 2

\pages 280--285



\Bibitem{Kozhanov_12:en}

\paper On a Nonlocal Boundary Value Problem with Variable Coefficients for the Heat Equation and the Aller Equation

\by  Kozhanov~A.~I.

\mathscinet{http://www.ams.org/mathscinet-getitem?mr=2162443}

\jour Differ. Equ.

\yr 2004

\vol 40

\issue 6

\pages 815--826

\crossref{10.1023/B:DIEQ.0000046860.84156.f0}



\Bibitem{Nakhushev_13:en}

\lang In Russian

\by Nakhushev~A.~M.

\book Uravneniia matematicheskoi biologii \rm [Equations of Mathematical Biology]

\yr 1995

\publ Vysshaia shkola

\publaddr Moscow

\totalpages 301





\Bibitem{Repin_14:en}

\paper An Analog of the Nakhushev Problem for the Bitsadze–Lykov Equation

\by Repin~O.~A.

\jour Differ. Equ.

\yr 2002

\vol 38

\issue 10

\pages 1503--1509

\crossref{10.1023/A:1022339217281}

\mathscinet{http://www.ams.org/mathscinet-getitem?mr=2014239}









\Bibitem{Repin_15:en}

\paper A nonlocal problem for the Bitsadze--Lykov equation

\by Repin~O.~A.,  Kumykova~S.~K.

\jour Russian Math. (Iz. VUZ)

\yr 2010

\vol 54

\issue 3

\pages 24--30

\crossref{10.3103/S1066369X10030059}

\scopus{http://www.scopus.com/record/display.url?origin=inward&eid=2-s2.0-78649533261}





\Bibitem{Nakhushev_16:en}

\lang In Russian

\by Nakhushev~A.~M.

\book Drobnoe ischislenie i ego primenenie \rm [Fractional calculus and its application]

\yr 2003

\publ Fizmatlit

\publaddr Moscow

\totalpages 272











\Bibitem{Kalmenov_17:en}

\lang In Russian

\by Kal'menov~T.~Sh.

\paper A criterion for the uniqueness of the solution of the Darboux problem for a certain degenerate hyperbolic equation

\jour Differ. Uravn.

\yr 1971

\vol 7

\issue 1

\pages 178--181









\Bibitem{Kalmenov_18:en}

\lang In Russian

\by Kal'menov~T.~Sh.

\paper The Darboux problem for a certain degenerate equation

\jour Differ. Uravn.

\yr 1974

\vol 10

\issue 1

\pages 59--68





\Bibitem{Pulkina_19:en}

\paper Certain nonlocal problem for a degenerate hyperbolic equation

\by  Pul'kina~L.~S.

\jour Math. Notes

\yr 1992

\vol 51

\issue 3

\pages 286--290

\mathscinet{http://www.ams.org/mathscinet-getitem?mr=1172230}

\zmath{https://zbmath.org/?q=an:0812.35085}

\crossref{10.1007/BF01206393}



\Bibitem{Repin_20:en}

\lang In Russian

\by Repin~O.~A., Kumykova~S.~K.

\paper A~Boundary-value Problem with Shift for a~Hyperbolic Equation Degenerate in the Interior of a~Region

\jour Vestn. Samar. Gos. Tekhn. Univ., Ser. Fiz.-Mat. Nauki \rm [J. Samara State Tech. Univ., Ser. Phys. Math. Sci.]

\yr 2014

\issue 1(34)

\pages 37--47

\crossref{10.14498/vsgtu1280}





\Bibitem{Balkizov_21:}

\lang In Russian

\by Balkizov~Zh.~A.

\paper The first boundary value problem for a~degenerate hyperbolic equation

\jour Vladikavkaz. Mat. Zh.

\yr 2016

\vol 18

\issue 2

\pages 19--30

\mathnet{http://mi.mathnet.ru/vmj577}



\Bibitem{Balkizov_22:en}

\lang In Russian

\by Balkizov~Zh.~A.

\paper The boundary-value problem for a degenerate hyperbolic equation in the area

\jour Izvestiya Vuzov. Severo-Kavkazskii Region. Natural Science

\yr 2016

\issue 1

\pages 5--10

\crossref{10.18522/0321-3005-2016-1-5-10}









\Bibitem{Fikhtengolz_23:e}

\lang In Russian

\by Fikhtengol'ts~G.~M.

\book Kurs differentsial'nogo i integral'nogo ischisleniia \rm [Course of Differential

and Integral Calculus]

\bookvol 1

\yr 2003

\publ Fizmatlit

\publaddr Moscow

\totalpages 680







\Bibitem{Balkizov_24:en}

\paper Dirichlet boundary value problem for a third order parabolic-hyperbolic equation with degenerating type and order in the hyperbolicity domain

\by Balkizov~Zh.~A.

\jour Ufa Math. Journal

\yr 2017

\vol 9

\issue 2

\pages 25--39

\crossref{10.13108/2017-9-2-25}

\scopus{http://www.scopus.com/record/display.url?origin=inward&eid=2-s2.0-85023623624}







\Bibitem{Makaova_25:en}

\lang In Russian

\by Makaova R.Kh.

\paper The second boundary-value problem for generalized Hallaire equation with fractional Riemann--Liouville derivative

\jour Doklady Adygskoi (Cherkesskoi) Mezhdunarodnoi akademii nauk

\yr 2015

\vol 17

\issue 3

\pages 35--38















\Bibitem{Wright_26:en}

\by   Wright~E.~M.

\paper The generalized Bessel function of order greater than one

\jour Quart. J. Math. Oxford Ser.

\yr 1940

\vol os-11

\issue 1

\pages 36--48

\crossref{10.1093/qmath/os-11.1.36}







\Bibitem{Pskhu_27:en}

\by Pskhu~A.~V.

\paper Initial-value problem for a~linear ordinary differential equation of noninteger order

\jour Sb. Math.

\yr 2011

\vol 202

\issue 4

\pages 571--582

\zmath{https://zbmath.org/?q=an:1226.34005}

\adsnasa{http://adsabs.harvard.edu/cgi-bin/bib_query?2011SbMat.202..571P}

\crossref{10.1070/SM2011v202n04ABEH004156}

\isi{http://gateway.isiknowledge.com/gateway/Gateway.cgi?GWVersion=2&SrcApp=PARTNER_APP&SrcAuth=LinksAMR&DestLinkType=FullRecord&DestApp=ALL_WOS&KeyUT=000292829300005}

\scopus{http://www.scopus.com/record/display.url?origin=inward&eid=2-s2.0-79959815598}





\end{thebibliography}





%\end{document}